%% Chapitre 07 — Garde-fous, permissions et sécurité
%% RÈGLE : uniquement des commandes sémantiques bookforge.

\chapitre{Garde-fous, permissions et sécurité}

\introChapitre{
  Un agent qui peut écrire des fichiers et lancer des commandes peut aussi
  casser des choses. Ce chapitre mêle principes durables et réglages concrets :
  comment cadrer ce que l'agent a le droit de faire, isoler les opérations
  risquées, et garder une marche arrière. La confiance se construit sur des
  garde-fous, pas sur de l'espoir.
}

\section{Le modèle de permissions : ce qui passe, ce qui demande approbation}

\paragraphe{Au chapitre précédent, tu as appris à relire et tester ce que
l'agent produit. C'est nécessaire, mais ça intervient \emphase{après} coup.
Mieux vaut prévenir : cadrer en amont ce que l'agent a le droit de faire. C'est
le rôle des \terme{garde-fou}s.}

\paragraphe{Le premier d'entre eux est intégré à Claude Code : le modèle de
\terme{permission}. Chaque action que l'agent veut entreprendre — lire un
fichier, en modifier un, lancer une commande — passe par une règle qui décide
de son sort. Trois issues possibles : l'action s'exécute directement, l'action
te demande une approbation, ou l'action est refusée.}

\paragraphe{Par défaut, le réglage est prudent. Lire un fichier ne coûte rien et
passe sans bruit. Mais dès que l'agent veut \emphase{écrire} ou exécuter une
commande shell, il s'arrête et te montre ce qu'il s'apprête à faire : tu vois le
\code{diff} ou la commande exacte, et tu approuves ou refuses.}

\paragraphe{Ce moment d'approbation est ta principale prise sur l'agent. Ne le
traite pas comme une formalité à expédier. Lis vraiment la commande proposée
avant de valider. L'agent qui veut lancer \code{npm test} est inoffensif ;
celui qui veut \code{rm -rf build/} mérite un coup d'œil attentif.}

\begin{astuce}
Quand l'agent te propose une action, Claude Code offre souvent une option du
genre « toujours autoriser ce type de commande ». Pratique, mais réfléchis :
tu ouvres une porte pour toute la session. Réserve-la aux actions vraiment
anodines et répétitives, comme lancer les tests.
\end{astuce}

\section{Commandes dangereuses et opérations destructrices}

\paragraphe{Certaines actions ne se rattrapent pas. Une suppression de fichiers,
un \code{git reset --hard}, un \code{git push --force}, une commande qui touche
une base de données ou un service distant : ce sont des opérations
\important{destructrices}. L'agent ne mesure pas toujours leur portée, parce
qu'il ne voit qu'une fraction de ton environnement dans son contexte.}

\paragraphe{Le risque n'est pas que l'agent soit malveillant — il ne l'est pas.
Le risque, c'est qu'il enchaîne une commande plausible mais inadaptée à \emphase{ta}
situation : effacer un dossier qu'il croit temporaire, écraser une branche qu'il
pense locale. Une hallucination sur l'état du système suffit.}

\begin{avertissement}
Ne mets jamais une commande destructrice en autorisation automatique. Les
\code{rm}, \code{git push --force}, \code{git reset --hard}, \code{DROP TABLE},
les déploiements et tout ce qui écrit hors du dépôt doivent rester en
approbation manuelle, sans exception. Le gain de temps d'un clic en moins ne
vaut pas une perte de données irréversible.
\end{avertissement}

\paragraphe{La parade durable est simple : tout ce qui est destructeur reste
sous ton contrôle explicite. Et quand une tâche \emphase{exige} une opération
risquée, fais-la toi-même plutôt que de la déléguer. L'agent prépare, tu
appuies sur le bouton rouge.}

\section{Travailler sur une branche : le filet du versionnage}

\paragraphe{Tu connais déjà Git. C'est ton meilleur garde-fou, et il ne coûte
rien. Avant de lancer l'agent sur une tâche, crée une branche dédiée :}

\begin{extraitcode}{bash}
git checkout -b agent/refacto-paiement
# ... tu laisses l'agent travailler sur cette branche ...
git diff main          # tu compares tout ce qu'il a changé
git checkout main      # marche arrière instantanée si ça part en vrille
git branch -D agent/refacto-paiement
\end{extraitcode}

\paragraphe{Sur une branche, l'agent peut se tromper sans conséquence durable.
Tout son travail est isolé du code stable. Si le résultat te convient, tu fais
une revue de \code{diff} puis tu fusionnes. Sinon, tu jettes la branche et il ne
reste aucune trace. La marche arrière est gratuite et totale.}

\paragraphe{Commite par petits pas pendant la session. Des commits fréquents
transforment chaque étape en point de sauvegarde : si la troisième modification
casse tout, tu reviens à la deuxième sans perdre les deux premières. Le
versionnage fin est un garde-fou plus puissant que n'importe quel réglage de
permission.}

\begin{astuce}
N'oublie pas le travail non commité. Avant une session ambitieuse, lance un
\code{git stash} ou un commit de tes changements en cours. L'agent ne respecte
pas tes modifications non sauvegardées : il peut les écraser sans le savoir.
\end{astuce}

\section{Isoler l'agent : répertoires de travail et bacs à sable}

\paragraphe{Le deuxième principe durable, c'est l'isolation : restreindre le
périmètre physique dans lequel l'agent opère. Claude Code travaille dans le
répertoire où tu le lances et ses sous-dossiers. Lance-le donc \emphase{depuis
le dépôt concerné}, jamais depuis ta racine personnelle ou un dossier qui
contient des projets sans rapport.}

\paragraphe{Plus le périmètre est large, plus une erreur peut faire de dégâts.
Un agent ouvert sur ton dossier personnel pourrait, sur un malentendu, toucher
des fichiers qui n'ont rien à voir avec la tâche. Réduis sa surface d'action à
ce dont il a strictement besoin.}

\paragraphe{Pour les tâches vraiment risquées — du code que tu n'as pas écrit,
une dépendance douteuse, une expérimentation — pousse l'isolation plus loin :
fais tourner l'agent dans un conteneur ou une machine jetable. Un bac à sable
te garantit que même une commande catastrophique reste enfermée et n'atteint ni
tes autres projets, ni ton système.}

\section{Secrets et données sensibles : ce qu'il ne faut jamais exposer}

\paragraphe{L'agent lit ton dépôt pour se constituer son contexte. S'il y trouve
des mots de passe, des clés d'API ou des jetons d'accès, ils entrent dans la
conversation — et donc potentiellement dans des journaux ou un \code{diff} que
tu partages plus tard.}

\begin{avertissement}
Ne laisse jamais traîner de secrets en clair dans ton dépôt. Clés d'API,
identifiants de base de données, jetons : ils vont dans un fichier
\code{.env} ignoré par Git, jamais dans le code. Ajoute ce fichier à ton
\code{.gitignore} pour que l'agent — comme Git — ne le versionne pas par
accident.
\end{avertissement}

\paragraphe{Tu peux aussi soustraire des chemins entiers à la lecture de
l'agent, même s'ils sont présents dans le dépôt. Le mécanisme courant passe par
les règles de permission, dans \code{.claude/settings.json} :}

\begin{extraitcode}{json}
{
  "permissions": {
    "deny": [
      "Read(./.env)",
      "Read(./.env.*)",
      "Read(./secrets/**)",
      "Read(./**/*.pem)",
      "Read(./deploy/credentials.json)"
    ]
  }
}
\end{extraitcode}

\paragraphe{Cette configuration ne se contente pas de cacher les fichiers dans
une liste : elle refuse explicitement les lectures qui correspondent aux motifs
déclarés. C'est plus robuste qu'un simple oubli volontaire, parce que la règle
s'applique même si l'agent tombe sur le chemin pendant son exploration.}

\paragraphe{Pense aussi aux secrets qui ne sont pas dans le dépôt mais dans ta
session : ne colle pas une clé de production dans le chat « pour aller plus
vite ». Et méfie-toi des données réelles d'utilisateurs : pour les tests,
préfère des jeux de données anonymisés.}

\section{Définir tes propres règles d'autorisation et de refus}

\paragraphe{Les réglages par défaut sont prudents, mais génériques. Tu peux les
ajuster à ton projet via un fichier de configuration : une liste d'actions
\emphase{toujours autorisées} (pour fluidifier le quotidien) et une liste
d'actions \emphase{toujours refusées} (pour fermer définitivement des portes).}

\begin{extraitcode}{json}
{
  "permissions": {
    "allow": [
      "Bash(npm test)",
      "Bash(git status)",
      "Read(src/**)"
    ],
    "deny": [
      "Bash(rm *)",
      "Bash(git push --force*)",
      "Read(./.env)",
      "Read(./.env.*)",
      "Read(./secrets/**)"
    ]
  }
}
\end{extraitcode}

\paragraphe{La règle d'or : sois généreux sur l'\emphase{allow} pour les actions
lisibles et réversibles, strict sur le \emphase{deny} pour tout ce qui est
destructeur ou sensible. Une bonne liste de refus agit comme une ceinture de
sécurité : tu ne la sens pas tant que tout va bien, et elle te sauve le jour où
l'agent dérape.}

\begin{avertissement}
La syntaxe exacte de ce fichier, son emplacement et le nom des règles évoluent
au fil des versions de Claude Code. Vérifie la documentation de ta version
plutôt que de recopier un exemple daté. Le \emphase{principe} — autoriser
l'anodin, refuser le dangereux — reste valable ; les détails, eux, peuvent
bouger.
\end{avertissement}

\resumeChapitre{
  \element{\important{Permissions} : lis vraiment chaque action proposée avant d'approuver — la demande de confirmation est ta principale prise sur l'agent.}
  \element{\important{Opérations destructrices} : \code{rm}, \code{git push --force}, \code{DROP TABLE} et tout déploiement restent en approbation manuelle sans exception.}
  \element{\important{Branche dédiée} : travaille sur une branche avant de lancer l'agent — si ça part en vrille, \code{git checkout main} ramène tout à zéro en une commande.}
  \element{\important{Isolation} : lance l'agent depuis le dépôt concerné uniquement, jamais depuis ta racine personnelle, pour réduire sa surface d'action au strict nécessaire.}
  \element{\important{Secrets} : les clés et identifiants vont dans \code{.env} ignoré par Git ; ajoute des règles \code{permissions.deny} pour refuser explicitement la lecture des fichiers sensibles.}
  \element{\important{Règles sur mesure} : liste blanche pour les actions anodines et répétitives, liste noire pour le destructeur et le sensible — la ceinture qui ne se fait sentir que quand ça compte.}
}

\paragraphe{Tu disposes maintenant d'un jeu de garde-fous complet : permissions,
branches, isolation, protection des secrets, règles sur mesure. Ces bases
serviront partout dans la suite — quand tu automatiseras des actions, brancheras
des outils externes ou affronteras les pièges classiques. Encadrer l'agent,
c'est ce qui rend la délégation sereine.}
